\section{CIS/TDA-Type Wave Functions}\label{sec-diabat-wfn-cis-tda}
For CIS and TDA, an electronic state is expressed as a linear combination of a reference determinant and its single excitations,
\begin{equation}
|\Psi_I\rangle=c_{\mathrm{ref}}^{(I)}|\Phi_0\rangle+
\sum_{\sigma=\alpha,\beta}\sum_{ia}X_{ia\sigma}^{(I)}
|\Phi_{i_\sigma}^{a_\sigma}\rangle.
\end{equation}
Here, $i$ and $a$ denote occupied and virtual molecular orbitals (MOs), and $\sigma$ denotes the electron spin. The coefficient $c_{\mathrm{ref}}^{(I)}$ describes the reference configuration, while $X_{ia\sigma}^{(I)}$ gives the amplitude for exciting an electron from orbital $i$ to orbital $a$. The same file format is used for TDA-aTB wave functions.

\subsection*{Complete CIS/TDA example}
Consider a closed-shell two-electron system with one occupied and one virtual spatial orbital. The following file describes a single TDA excited state with one nonzero excitation amplitude.
\begin{lstlisting}[style=technicalfile]
## Wave-function data exported by Diabat
Method: TDA
Orbitals are not spin-polarized.
Number of electrons: 2
Number of alpha electrons: 1
Number of beta electrons: 1
Number of occupied orbitals (alpha): 1
Number of virtual orbitals (alpha): 1
Number of occupied orbitals (beta): 1
Number of virtual orbitals (beta): 1
Total energy of the reference ground state (a.u.): -1.0

State Information:

  State index: 1

  Total energy (a.u.): -0.9
  Excitation energy (eV): 2.7211386
  <S^2>: 0.0
  Amplitudes:
  ------------------------------------------------------------
       1   -->      2    7.071067811865475E-01  alpha
  ------------------------------------------------------------
\end{lstlisting}
The example uses spin-unpolarized orbitals, so the listed alpha amplitude represents both spin channels. Its value of $1/\sqrt{2}$ gives a normalized state under the convention described below. The energies and orbital counts are illustrative and must be replaced with values consistent with the accompanying MO file.

\subsection*{Excitation amplitudes}
The file header specifies whether the orbitals are spin-polarized, the numbers of alpha and beta electrons, the numbers of occupied and virtual orbitals for each spin, and the total energy of the reference ground state. Use \texttt{Orbitals are not spin-polarized.} when both spins share the same spatial orbitals. Otherwise use \texttt{Orbitals are spin-polarized.} and provide the alpha and beta excitation amplitudes separately.

Each coefficient record contains an occupied orbital index, \texttt{-->}, a virtual orbital index, a real amplitude, and the spin label \texttt{alpha} or \texttt{beta}. Orbital indices start at 1 and follow the accompanying MO file. Virtual orbitals retain their full MO indices rather than being numbered again from 1. In the example, \texttt{1 --> 2} denotes excitation from the first occupied orbital to the first virtual orbital. Only \texttt{-->} records are used for CIS and TDA.

If a state contains a reference-configuration contribution, write its coefficient as
\begin{lstlisting}[style=technicalfile]
        Ref. config.    2.000000000000000E-02
\end{lstlisting}
This record belongs inside the \texttt{Amplitudes:} block. An omitted reference coefficient or excitation amplitude is taken as zero.

\subsection*{Normalization}
The coefficient norm is
\begin{equation}
\mathcal{N}=|c_{\mathrm{ref}}|^2+\sum_{\sigma=\alpha,\beta}\sum_{ia}|X_{ia\sigma}|^2.
\end{equation}
For spin-unpolarized orbitals, \progname{} stores only alpha amplitudes and counts its contribution twice. For spin-polarized orbitals, the alpha and beta amplitudes are stored and counted separately. In a \texttt{statepack} block, the reference ground state may be loaded using \texttt{state\_index = 0} without including a separate state record in the file.
